\section*{Capítulo \ref{ch:b3:extension-strike-slip} --- Tectónica extensional y de desgarre}

\begin{solution}{exo:b3:extension-strike-slip:1}
Los estratos buzan $50 - 35 = \ang{15}$, alejándose de las fallas; $\beta =
\sin\ang{50}/\sin\ang{35} = 0.766/0.574 = 1.34$.
\end{solution}

\begin{solution}{exo:b3:extension-strike-slip:2}
Escalonamiento dextral hacia la derecha: cuenca. Escalonamiento dextral hacia la izquierda: dorsal. Escalonamiento sinistral
hacia la izquierda: cuenca. Escalonamiento sinistral hacia la derecha: dorsal.
\end{solution}

\begin{solution}{exo:b3:extension-strike-slip:3}
$105/20 = \qty{5.25}{km/Myr} = \qty{5.25}{mm/yr}$, cercana a la
\qty{6}{mm/yr}: la tasa no tiene por qué haber sido constante durante \qty{20}{Myr}.
\end{solution}

\begin{solution}{exo:b3:extension-strike-slip:4}
\qty{200}{\percent} significa $\Delta L = 2L_0$: $L_0 = \qty{117.5}{km}$ y $L
= \qty{352.5}{km}$. En el norte, $\Delta L = 0.5L_0$: $L_0 = \qty{470}{km}$
y $L = \qty{705}{km}$.
\end{solution}

\begin{solution}{exo:b3:extension-strike-slip:5}
De $563/20 = 28$ a $563/15 = \qty{38}{mm/yr}$. El desplazamiento es un mínimo, y
la deriva medida mediante levantamientos geodésicos a través de la zona de falla también incluye el movimiento en
fallas vecinas del sistema y la deformación elástica acumulada entre
terremotos.
\end{solution}

\begin{solution}{exo:b3:extension-strike-slip:6}
$t = 90 - \text{buzamiento}$. Para \ang{60}: $(1 + 0.6\times1.732)/(1 - 0.6\times
0.577) = 2.039/0.654 = 3.12$; para \ang{50}: $1.715/0.497 = 3.45$; para
\ang{40}: $1.504/0.285 = 5.28$. Bloqueo a $\arctan0.6 = \ang{31.0}$.
\end{solution}

\begin{solution}{exo:b3:extension-strike-slip:7}
$\arctan0.4 = \ang{21.8}$ y $\arctan0.2 = \ang{11.3}$. Una falla que buza
\ang{15} solo desliza si $\mu < \tan\ang{15} = 0.27$.
\end{solution}

\begin{solution}{exo:b3:extension-strike-slip:8}
$2.5\times4.0 = \qty{10.0}{km^2}$ y $2.5\times(4.0 - 1.5) =
\qty{6.25}{km^2}$.
\end{solution}

\begin{solution}{exo:b3:extension-strike-slip:9}
$132/1.19 = \qty{110.9}{km}$, con un alargamiento de \qty{21.1}{km}, lo que requiere
$21.1/6 = \qty{3.5}{Myr}$.
\end{solution}

\begin{solution}{exo:b3:extension-strike-slip:10}
$3.0/\sin\ang{60} = \qty{3.464}{km}$ antes y $3.0/\sin\ang{45} =
\qty{4.243}{km}$ después: \qty{0.779}{km} por bloque ($\beta = 1.22$). Los estratos
buzan $60 - 45 = \ang{15}$. Cada bloque gira como un todo: la rotación que
reduce el buzamiento de las fallas que buzan, por ejemplo, hacia el E, gira una línea horizontal del
bloque de modo que se eleve hacia el E, y los estratos buzan hacia el O, en sentido contrario a las
fallas.
\end{solution}

\begin{solution}{exo:b3:extension-strike-slip:11}
En una zona sinistral E--O $\sigma_1$ se orienta a \ang{045}. Las cizallas R a
$\phi/2 = \ang{20}$ de la zona, en el lado de $\sigma_1$: rumbo \ang{070},
sinistral. Las R$'$ a $90 - 20 = \ang{70}$: rumbo \ang{020}, dextral. Las fracturas de tensión
paralelas a $\sigma_1$: rumbo \ang{045}.
\end{solution}

\begin{solution}{exo:b3:extension-strike-slip:12}
Para $\mu = 0.6$ las fallas se forman a \ang{60.5} y se bloquean a \ang{31.0}: $\beta =
\sin\ang{60.5}/\sin\ang{31.0} = 0.870/0.515 = 1.69$. Para $\mu = 0.85$: de
\ang{65.2} a \ang{40.4}, $\beta = 0.908/0.648 = 1.40$. Las fallas con menor
fricción rotan más antes de bloquearse: las rocas débiles se extienden más antes de que
deban formarse nuevas fallas.
\end{solution}

\begin{solution}{pb:b3:extension-strike-slip:1}
\textbf{1.} Los estratos estaban horizontales y han girado con los bloques; las
fallas han girado el mismo ángulo, por lo que $\theta_0 = 42 + 18 = \ang{60}$,
el buzamiento de Anderson para $\phi = \ang{30}$.

\textbf{2.} $\beta = \sin\ang{60}/\sin\ang{42} = 0.8660/0.6691 = 1.294$.

\textbf{3.} $24.0/1.294 = \qty{18.5}{km}$; el dominio se ha alargado en
\qty{5.5}{km} (\qty{5.46}{km}).

\textbf{4.} $t = \ang{48}$: $(1 + 0.6\times0.900)/(1 - 0.6\times1.111) =
1.540/0.334 = 4.62$, frente a $3.12$ de las fallas nuevas a \ang{60.5}: casi
la mitad más.

\textbf{5.} En $\arctan0.6 = \ang{31.0}$; $\beta = \sin\ang{60}/\sin\ang{31}
= 1.68$ como máximo. Antes de eso, nuevas fallas a unos \ang{60} cortarían los
bloques, y una segunda generación de \omterm{def:b3:extension-strike-slip:domino}{fallas en dominó} tomaría el relevo.

\textbf{6.} Cada columna del bloque superior desciende cuanto desciende la falla situada
bajo ella a lo largo del salto horizontal; la falla es más empinada cerca de la superficie y se aplana
a medida que se aleja de ella, de modo que las columnas próximas a la falla descienden más y las más alejadas
descienden menos: los estratos se inclinan hacia la falla.

\textbf{7.} $D = A/H = 16.0/3.2 = \qty{5.0}{km}$.

\textbf{8.} $D$ es proporcional a $A$: $\qty{5.0}{km} \pm \qty{0.5}{km}$.

\textbf{9.} $3.2\times(5.0 - 2.0) = \qty{9.6}{km^2}$.

\textbf{10.} El salto horizontal, \qty{3.2}{km}.

\textbf{11.} El relleno se engrosa hacia la falla, en forma de cuña, y sus capas más antiguas
están más inclinadas que las más jóvenes: cada capa se depositó mientras
el \omterm{def:b3:extension-strike-slip:detachment}{anticlinal de rollover} y el espacio junto a la falla seguían creciendo.

\textbf{12.} $z = P/(\rho g) = \num{5.5e8}/(2800\times9.8) = \qty{20.0}{km}$.

\textbf{13.} $20.0/(24 - 16) = \qty{2.5}{km/Myr} = \qty{2.5}{mm/yr}$.

\textbf{14.} $20.0/\sin\ang{30} = \qty{40}{km}$ de desplazamiento, y $40\cos\ang{30} =
\qty{34.6}{km}$ de extensión horizontal.

\textbf{15.} No: \ang{15} está por debajo del buzamiento de bloqueo de \ang{31}. O bien la
falla se desplazó mientras era más empinada y luego se rotó (\omterm{def:b3:extension-strike-slip:core-complex}{charnela rodante}), o bien
era débil, con $\mu < 0.27$ debido a arcillas o fluidos a alta presión.

\textbf{16.} Por $45 - 15 = \ang{30}$: al alejarse la placa superior, el
bloque inferior descargado se elevó por isostasia y se curvó, rotando la falla a través de la
charnela.

\textbf{17.} \qty{34.6}{km} frente a los \qty{8.7}{km} del semigraben (pregunta
23): cuatro veces más. Los \omterm{def:b3:extension-strike-slip:core-complex}{complejos metamórficos de núcleo} se sitúan donde la extensión es mayor, como en
la cuenca y cordillera central (Basin and Range).

\textbf{18.} Una falla sinistral que da un escalón hacia la izquierda: distensiva, se abre una cuenca.

\textbf{19.} $\qty{5}{mm/yr}\times\qty{1.2}{Myr} = \qty{6.0}{km}$.

\textbf{20.} $4.0\times6.0 = \qty{24}{km^2}$; $3.0 + 6.0 = \qty{9.0}{km}$ de longitud;
$9.0/4.0 = 2.25$.

\textbf{21.} $132/10 = 13$ a $132/7 = 19$: la cuenca del mar Muerto es mucho más
alargada, con un desplazamiento muchas veces mayor que su anchura.

\textbf{22.} Una \omterm{def:b3:extension-strike-slip:bends}{curvatura contractiva}: una cresta empujada hacia arriba entre los ramales; en
sección, una \omterm{def:b3:extension-strike-slip:flower}{estructura en flor} positiva, con la falla dividida hacia arriba en ramales
con componente inversa, como la de entre 100 y \qty{300}{m} de anchura en la
falla transformante del mar Muerto.

\textbf{23.} $3.2 + 5.46 = \qty{8.66}{km}$: el semigraben de Merrow se ha
extendido \textbf{\qty[text-series-to-math=true]{8.7}{km}}.
\end{solution}
